\documentclass[10pt,a4paper]{amsart}
\usepackage[margin=19mm]{geometry}
\usepackage{amsmath,amssymb,mathtools}
\usepackage[scr=rsfs,cal=euler]{mathalfa}
\usepackage{xfrac}
\usepackage[unicode,hidelinks,bookmarks=false]{hyperref}
\newcommand{\ie}{i.e.\ }
\newcommand{\cf}{cf.\ }
\newcommand{\resp}{respectively\ }
\newcommand{\Subsection}{\subsection}
\providecommand{\nobreakdash}{\nobreak\hbox{-}}
\setlength{\emergencystretch}{2em}
\multlinegap0pt
\usepackage{mathtools}
\usepackage{amssymb}
\usepackage{enumerate}
\newtheorem{theorem}{Theorem}
\title{Comparison Results for a class of Neumann Problems of the $p$-Laplace Equation on Riemannian Manifolds}
\author{Wentao Liu, Anqiang Zhu}
\date{Source: arXiv:2606.22854v1 (CC BY 4.0)}
\begin{document}
\maketitle

\noindent\emph{Source and licence.} Comparison Results for a class of Neumann Problems of the $p$-Laplace Equation on Riemannian Manifolds. Version arXiv:2606.22854v1 (CC BY 4.0), distributed by arXiv under the Creative Commons Attribution 4.0 International licence, http://creativecommons.org/licenses/by/4.0/, as recorded in the arXiv licence metadata for this exact version and shown on the abstract page https://arxiv.org/abs/2606.22854v1. Full licence notice: \texttt{licenses/arxiv-2606.22854-CC-BY-4.0.txt}.\par\smallskip

\noindent\textit{Editorial scope.} Counted unit: the article's theorem thm:1.3 (source0.tex lines 309--349). The article's complete original proof of it is delivered with it (source0.tex lines 1145--1383, one \texttt{proof} environment of 239 lines, which the article places away from the statement). No mathematics is written, repaired or re-derived: the statement and the proof are the article's own lines, in the article's own notation, and no retained line is substituted or re-flowed.  Retained uncounted as the source-local context the proof needs: lines 151--160; lines 166--181; lines 247--251; lines 269--278; lines 280--289; lines 302--307; lines 627--640; lines 775--779.  Nothing was omitted from any retained span, so the package carries no omission record.  No retained line contains a citation, so the wrapper carries no bibliography and no bibliography entry.  No retained line contains a figure environment, an image-inclusion command or drawing code, so no image is delivered and the package declares no asset dependency.  Editorial formatting only, in addition: an AMS \texttt{amsart} preamble that restates the article's own declarations and macros used by the retained lines, namely the environments \texttt{theorem} on separate counters, together with the standard packages those lines need.  The retained environments print in this document as Theorem 1 in order of appearance, whereas the article prints the same items with the numbering of its own full document.  The complete native source of the article as distributed in the arXiv e-print for this exact version is bundled unchanged as \texttt{sources/203387/source0.tex}.
\begin{equation}\label{eq:1.3}
\begin{cases}
-\operatorname{div}\bigl(|\nabla u|^{p-2}\nabla u\bigr)=f,
& \text{in }\Omega,
\\[1ex]
|\nabla u|^{p-2}\dfrac{\partial u}{\partial\nu}
+\beta^{p-1}=0,
& \text{on }\partial\Omega.
\end{cases}
\end{equation}

\begin{equation}\label{eq:1.4}
\begin{cases}
-\operatorname{div}\bigl(|\nabla v|^{p-2}\nabla v\bigr)
=
f^\sharp,
& \text{in }\Omega^\sharp,
\\[1ex]
|\nabla v|^{p-2}
\dfrac{\partial v}{\partial\nu}
+
(\beta^\sharp)^{p-1}
=
0,
& \text{on }\partial\Omega^\sharp.
\end{cases}
\end{equation}

\begin{align}\label{eq:1.7}
\frac{\beta^\sharp}{\beta}\,\alpha_q
&=
\theta_\kappa\,\alpha_q^\sharp.
\end{align}

\begin{align}\label{eq:1.8}
\|u\|_{L^{\lambda,1}(\Omega)}
\le
\theta_\kappa^{\frac{1}{\lambda}}
\|v\|_{L^{\lambda,1}(\Omega^\sharp)},
\qquad
\kappa=0,1,
\quad
0 < \lambda \le \frac{n(p-1)}{p(n-1)},
\end{align}

\begin{align}\label{eq:1.9}
\|u\|_{L^{q\lambda,q}(\Omega)}
\le
\theta_\kappa^{\frac{1}{q\lambda}}
\|v\|_{L^{q\lambda,q}(\Omega^\sharp)},
\qquad
\kappa=0,
\quad
0 < \lambda \le \frac{n(p-1)}{p(n-2)+n},
\end{align}

\begin{align}\label{eq:1.11}
\|u\|_{L^{q\lambda,q}(\Omega)}
\le
\theta_\kappa^{\frac{1}{q\lambda}}
\|v\|_{L^{q\lambda,q}(\Omega^\sharp)}.
\end{align}

\begin{align}
S_\kappa(\mu(t))^{\frac{p}{p-1}}
&\le
\left(\int_0^{\mu(t)} f^*(s)\,ds\right)^\frac{1}{p-1}
\left(-\mu'(t)+\frac{1}{\beta}\int_{\partial U_t^e} d\sigma\right),
\label{eq:3.1}
\\
\tilde S_\kappa(\varphi(t))^\frac{p}{p-1}
&=
\theta_\kappa^{-\frac{1}{p-1}}
\left(\int_0^{\theta_\kappa \varphi(t)} f^*(s)\,ds\right)^{\frac{1}{p-1}}
\left(-\varphi'(t)+\frac{1}{\beta^\sharp}\int_{\partial V_t^e} d\sigma\right),
\label{eq:3.2}
\end{align}

\begin{align}\label{eq:3.3}
\mu(t) \le \theta_\kappa \varphi(t),
\qquad
\text{for all } 0\le t \le v_m.
\end{align}

\begin{theorem}\label{thm:1.3}
Assume that $f\equiv1$ in $\Omega$.
Let $q\ge 1$, and let $u$ be a positive weak solution to problem \eqref{eq:1.3}.

Suppose that $v$ is the positive solution to problem \eqref{eq:1.4}
satisfying condition \eqref{eq:1.7}. If
\[
1\le p\le \frac{n}{n-1},
\]
then
\begin{align}
u^\sharp(x)\le \theta_\kappa v(x),
\qquad
\forall\,x\in\Omega^\sharp.
\end{align}

If
\[
p > \frac{n}{n-1},
\]
then the following estimates hold:
\begin{align}\label{eq:1.13}
\|u\|_{L^{\lambda,1}(\Omega)}
\le
\theta_\kappa^{\frac{1}{\lambda}}
\|v\|_{L^{\lambda,1}(\Omega^\sharp)},
\qquad
\kappa=0,1,
\quad
0<\lambda \le\frac{n(p-1)}{n(p-1)-p}.
\end{align}
Moreover, when $\kappa=0$, one has
\begin{align}\label{eq:1.14}
\|u\|_{L^{q\lambda,q}(\Omega)}
\le
\theta_\kappa^{\frac{1}{q\lambda}}
\|v\|_{L^{q\lambda,q}(\Omega^\sharp)},
\qquad
0<\lambda \le\frac{n(p-1)}{n(p-1)-p}.
\end{align}
\end{theorem}

\begin{proof}[Proof of Theorem \ref{thm:1.3}]
  For the special case $f\equiv 1$, we first observe that
\[
\int_0^{\mu(t)} f\, ds = \mu(t), 
\qquad 
\int_0^{\theta_\kappa \varphi(t)} f\, ds = \theta_\kappa \varphi(t).
\]
Hence, equations \eqref{eq:3.2} and \eqref{eq:3.3} reduce to
\begin{align}
S_\kappa(\mu(t))^{\frac{p}{p-1}}
&\le
\mu(t)^{\frac{1}{p-1}}
\left(-\mu'(t)+\frac{1}{\beta}\int_{\partial U_t^e} d\sigma\right),
\label{eq:3.7}
\\
\tilde S_\kappa(\varphi(t))^{\frac{p}{p-1}}
&=
\varphi(t)^{\frac{1}{p-1}}
\left(-\varphi'(t)+\frac{1}{\beta^\sharp}\int_{\partial V_t^e} d\sigma\right),
\label{eq:3.8}
\end{align}
For \eqref{eq:3.7}, we multiply both sides by
$t^{q-1} S_\kappa(\mu(t))^{-\frac{p}{p-1}}$
and integrate over $(0,\tau)$, with $\tau \ge v_m$. This yields
\begin{align*}
\int_0^\tau t^{q-1}\,dt
&\le
\int_0^\tau (-\mu'(t))\, t^{q-1}\mu(t)^{\frac{1}{p-1}}
S_\kappa(\mu(t))^{-\frac{p}{p-1}}\,dt
\nonumber\\
&\quad +
\frac{1}{\beta}
\int_0^\tau t^{q-1}\mu(t)^{\frac{1}{p-1}}
S_\kappa(\mu(t))^{-\frac{p}{p-1}}
\int_{\partial U_t^e} d\sigma\, dt.
\end{align*}

Since the function
$l \mapsto l^{\frac{1}{p-1}} S_\kappa(l)^{-\frac{p}{p-1}}$
is nondecreasing for $1 \le p \le \frac{n}{n-1}$, we deduce that
\begin{align*}
\int_0^\tau t^{q-1}\,dt
\le
\int_0^\tau (-\mu'(t))\, t^{q-1}\mu(t)^{\frac{1}{p-1}}
S_\kappa(\mu(t))^{-\frac{p}{p-1}}dt
+
\frac{\alpha_q}{\beta}
|\Omega|^{\frac{1}{p-1}}
S_\kappa(|\Omega|)^{-\frac{p}{p-1}}.
\end{align*}
We apply the same argument to \eqref{eq:3.8}. Consequently, we obtain
\begin{align*}
\int_0^\tau t^{q-1}\,dt
=
\int_0^\tau (-\varphi'(t))\, t^{q-1}\varphi(t)^{\frac{1}{p-1}}
\tilde{S}_\kappa(\varphi(t))^{-\frac{p}{p-1}}\,dt
+\theta_\kappa\frac{\alpha_q^\sharp}{\beta^\sharp}
|\Omega|^{\frac{1}{p-1}}
S_\kappa(|\Omega|)^{-\frac{p}{p-1}}.
\end{align*}

Combining the above inequality with the corresponding estimate for $\mu(t)$, we deduce that
\begin{align*}
\int_0^\tau (-\mu'(t))\,t^{q-1}\mu(t)^{\frac{1}{p-1}}
S_\kappa(\mu(t))^{-\frac{p}{p-1}}\,dt
\ge
\int_0^\tau (-\varphi'(t))\,t^{q-1}\varphi(t)^{\frac{1}{p-1}}
\tilde{S}_\kappa(\varphi(t))^{-\frac{p}{p-1}}\,dt.
\end{align*}

An integration by parts yields
\begin{align*}
&-\tau^{p-1}\int_0^{\mu(\tau)} l^{\frac{1}{p-1}} S_\kappa(l)^{-\frac{1}{p-1}}\,dl
+ (p-1)\int_0^\tau t^{p-2}\int_0^{\mu(t)} l^{\frac{1}{p-1}} S_\kappa(l)^{-\frac{p}{p-1}}\,dl\,dt \\
&\ge
-\tau^{p-1}\int_0^{\theta_\kappa \varphi(\tau)} l^{\frac{1}{p-1}} S_\kappa(l)^{-\frac{1}{p-1}}\,dl
+ (p-1)\int_0^\tau t^{p-2}\int_0^{\theta_\kappa \varphi(t)} l^{\frac{1}{p-1}} {S}_\kappa(l)^{-\frac{p}{p-1}}\,dl\,dt.
\end{align*}
We apply Gronwall's inequality once again, with
\[
\zeta(\tau)
=
\int_0^\tau t^{p-2}
\int_{\theta_\kappa \varphi(t)}^{\mu(t)}
l^{\frac{1}{p-1}} S_\kappa(l)^{-\frac{p}{p-1}}\,dl\,dt,
\qquad C=0.
\]
It follows that
\begin{align*}
\tau^{p-2}
\int_{\theta_\kappa \varphi(\tau)}^{\mu(\tau)}
l^{\frac{1}{p-1}} S_\kappa(l)^{-\frac{p}{p-1}}\,dl
\le
\frac{(q-1)\zeta(v_m)}{v_m}
\left(\frac{\tau}{v_m}\right)^{q-2}.
\end{align*}

Since $\zeta(v_m)\le 0$ and the above inequality holds for every $\tau \ge v_m$, we deduce that
\[
\int_{\theta_\kappa \varphi(\tau)}^{\mu(\tau)}
l^{\frac{1}{p-1}} S_\kappa(l)^{-\frac{p}{p-1}}\,dl
\le 0,
\qquad \forall\, \tau \ge v_m.
\]

On the other hand, since
$l^{\frac{1}{p-1}} S_\kappa(l)^{-\frac{p}{p-1}} > 0$
in the integration range, it follows that
\[
\mu(\tau) \le \theta_\kappa \varphi(\tau),
\qquad \forall\, \tau \ge v_m.
\]

Combining this with the previous estimate, we conclude that
\[
\mu(t) \le \theta_\kappa \varphi(t), \qquad \forall\, t>0,
\]
and therefore \eqref{eq:1.11} holds.
\par
We now turn to the proof of \eqref{eq:1.13} and \eqref{eq:1.14}. 
\par
Multiplying \eqref{eq:3.7} by
$t^{q-1} S_\kappa(\mu(t))^{-\frac{p}{p-1}}$
and integrating over $(0,\tau)$, with $\tau \ge v_m$, we obtain
\begin{align*}
\int_0^\tau t^{q-1}\,\mu(t)^{\frac{1}{\lambda}}\,dt
&\le
\int_0^\tau -t^{q-1}\,\mu'(t)\,\mu(t)^{\frac{1}{p-1}+\frac{1}{\lambda}}S_\kappa(\mu(t))^{-\frac{p}{p-1}}\,dt
\notag\\
&\quad
+\frac{1}{\beta}
\int_0^\tau t^{q-1}\,\mu(t)^{\frac{1}{p-1}+\frac{1}{\lambda}}S_\kappa(\mu(t))^{-\frac{p}{p-1}}
\left(\int_{\partial U_t^e} d\sigma\right) dt.
\end{align*}

Using the monotonicity of $l^{\frac{1}{p-1}+\frac{1}{\lambda}}S_\kappa(l)^{-\frac{p}{p-1}}$ for $\kappa=0,1$ when $0 < \lambda\le \frac{n(p-1)}{n(p-1)-p}$. We can obtain that
\begin{align}\label{eq:3.13}
\int_0^\tau t^{q-1}\,(\mu(t))^{\frac{1}{\lambda}}\,dt
\le
-\int_0^\tau t^{q-1}\,dF_\kappa((\mu(t)))
+\frac{\alpha_q}{\beta}
|\Omega|^{1+\frac{1}{\lambda}}S_\kappa(|\Omega|)^{-\frac{p}{p-1}}.
\end{align}
Proceeding as in the case of \eqref{eq:3.7}, we consider \eqref{eq:3.8}. Multiplying both sides by
\[
t^{q-1}\tilde{S}_\kappa(\varphi(t))^{-\frac{p}{p-1}}
\]
and integrating over $(0,\tau)$, $\tau \ge v_m$, we obtain
\begin{align}\label{eq:3.14}
\int_0^\tau t^{q-1}\,(\theta_\kappa\varphi(t))^{\frac{1}{\lambda}}\,dt
=
-\int_0^\tau t^{q-1}\,dF_\kappa((\theta_\kappa\varphi(t)))
+\frac{\alpha_q}{\beta}
|\Omega|^{1+\frac{1}{\lambda}}S_\kappa(|\Omega|)^{-\frac{p}{p-1}}.
\end{align}
Following the similar arguments in proof of \eqref{eq:1.8}, we are able to prove \eqref{eq:1.13}.
\par
We now turn to the proof of \eqref{eq:1.14}. It suffices to show that
\begin{align*}
\int_0^\infty t^{q-1}\,\mu(t)^{\frac{1}{\lambda}}\,dt
\le
\int_0^\infty t^{q-1}\,(\theta_0\varphi(t))^{\frac{1}{\lambda}}\,dt.
\end{align*}

To this end, Integrate \eqref{eq:3.13} and \eqref{eq:3.14} by parts on the right with the first term and letting $\tau\to\infty$. This reduces the problem to proving the equivalent inequality
\[
\int_0^\infty t^{q-1}F_0(\mu(t))\,dt
\le
\int_0^\infty t^{q-1}F_0(\theta_0 \varphi(t))\,dt.
\]
We multiply both sides of \eqref{eq:3.7} by
$t^{q-1} F_0(\mu(t)) S_0(\mu(t))^{\frac{1}{p-1}}$
and integrate over $(0,\tau)$, where $\tau \ge v_m$. This yields
\begin{align*}
\int_0^\tau t^{q-1} F_0(\mu(t))\,dt
&\le
\int_0^\tau -t^{q-1}\mu'(t)\, F_0(\mu(t)) S_0(\mu(t))^{-\frac{1}{p-1}}
\mu(t)^{\frac{1}{p-1}}\,dt \\
&\quad +
\frac{1}{\beta}
\int_0^\tau t^{q-1}\mu(t)^{\frac{1}{p-1}} F_0(\mu(t))
S_0(\mu(t))^{-\frac{p}{p-1}}
\int_{\partial U_t^e} d\sigma\, dt.
\end{align*}

Since the function
\[
l \mapsto l^{\frac{1}{p-1}} F_0(l) S_0(l)^{-\frac{p}{p-1}}
\]
is nondecreasing for
$0 < \lambda \le \frac{n(p-1)}{n(p-1)-p}$,
we deduce that
\begin{align*}
\int_0^\tau t^{q-1} F_0(\mu(t))\,dt
&\le
\int_0^\tau -t^{q-1}\mu'(t)\, F_0(\mu(t)) S_0(\mu(t))^{-\frac{1}{p-1}}
\mu(t)^{\frac{1}{p-1}}\,dt \\
&\quad +
\frac{\alpha_q}{\beta}
|\Omega|^{-\frac{1}{p-1}}
F_0(|\Omega|)
S_0(|\Omega|)^{-\frac{p}{p-1}}.
\end{align*}
We next apply a similar argument to \eqref{eq:3.8}. Multiplying both sides of \eqref{eq:3.8} By
\[
t^{q-1} F_0(\theta_0 \varphi(t)) \tilde{S}_0(\varphi(t))^{-\frac{p}{p-1}}
\]
and integrating over $(0,\tau)$, with $\tau \ge v_m$, we obtain
\begin{align*}
\int_0^\tau t^{q-1} F_0(\theta_0\varphi(t))\,dt
&=
\int_0^\tau -t^{q-1}\varphi'(t)\,F_0(\theta_0\varphi(t))
\tilde{S}_0(\varphi(t))^{-\frac{1}{p-1}}
\varphi(t)^{\frac{1}{p-1}}\,dt \\
&\quad +
\frac{1}{\beta^\sharp}
\int_0^\tau t^{q-1}\varphi(t)^{\frac{1}{p-1}}
F_0(\theta_0\varphi(t))
\tilde{S}_0(\varphi(t))^{-\frac{p}{p-1}}
\int_{\partial V_t^e} d\sigma\,dt.
\end{align*}

This can be rewritten as
\begin{align*}
\int_0^\tau t^{q-1} F_0(\theta_0\varphi(t))\,dt
&=
\int_0^\tau -t^{q-1}\varphi'(t)\,F_0(\theta_0\varphi(t))
\tilde{S}_0(\varphi(t))^{-\frac{1}{p-1}}
\varphi(t)^{\frac{1}{p-1}}\,dt \\
&\quad +
\theta_0
\frac{\alpha_q^\sharp}{\beta^\sharp}
|\Omega|^{-\frac{1}{p-1}}
F_0(|\Omega|)
S_0(|\Omega|)^{-\frac{p}{p-1}}.
\end{align*}

We now return to the proof of \eqref{eq:1.9}. Arguing in the same way as above, we conclude that \eqref{eq:1.14} holds.
\end{proof}

\end{document}
